\begin{trapbox}
	\begin{description}[style=nextline, leftmargin=2.8em, labelsep=0.5em, itemsep=0.35em]
		\item[\textbf{\textcolor{CTrap}{易错点一（仅翻折轴下方）：}}] 错误认为翻折就是将整个图像关于$x$轴对称，结果将$x$轴上方的部分也翻折下去，导致图像完全对称。
		\item[\textbf{\textcolor{CTrap}{正确理解（只翻折负值部分）：}}] 翻折只针对原函数值$f(x)<0$的部分，将它们关于$x$轴对称翻折到上方；对$f(x)\ge0$的部分，图像保持不变。
		\item[\textbf{\textcolor{CTrap}{错因分析（误用全局对称）：}}] 对绝对值定义理解不深刻，以为$|f(x)|$是$f(x)$的相反数，忽略了分段性。
	\end{description}

	\medskip
	\begin{description}[style=nextline, leftmargin=2.8em, labelsep=0.5em, itemsep=0.35em]
		\item[\textbf{\textcolor{CTrap}{易错点二（值域误判为无界）：}}] 误认为$y=|f(x)|$的值域一定是$[0,+\infty)$，而实际情况可能是有界区间。
		\item[\textbf{\textcolor{CTrap}{正确理解（值域取决原函数）：}}] 值域是原函数值域中每个值取绝对值后的集合，仍可能有上界，例如$f(x)=\sin x$，$|f(x)|=|\sin x|$的值域为$[0,1]$（但若原函数值域有界则可能有界）。
		\item[\textbf{\textcolor{CTrap}{错因分析（忽视原函数范围）：}}] 只注意到绝对值非负，却忽略了原函数本身的有界性，误以为翻折后必然无限向上延伸。
	\end{description}

	\medskip
	\begin{description}[style=nextline, leftmargin=2.8em, labelsep=0.5em, itemsep=0.35em]
		\item[\textbf{\textcolor{CTrap}{易错点三（对称轴混淆）：}}] 错误地将翻折理解为关于$y$轴对称（即左右翻折），导致图像画错。
		\item[\textbf{\textcolor{CTrap}{正确理解（关于x轴翻折）：}}] 翻折变换是关于$x$轴的对称，即纵坐标取相反数，横坐标不变；只有$x$轴下方的部分被翻折，$x$轴上方不变。
		\item[\textbf{\textcolor{CTrap}{错因分析（混淆坐标轴概念）：}}] 学习函数对称时，容易将关于$x$轴和关于$y$轴的对称记混，需要结合绝对值定义区分。
	\end{description}
\end{trapbox}
